% ============================================
% DOC-03 (EN)
% English translation — SPEC-935-R355
% ============================================
\section*{\textbf{DOC-03} --- Autopsy Report}
\markboth{DOC-03 --- Autopsy Report}{}

\noindent \textit{Forensic Medical Institute of Barbacena --- June 14, 1979 --- Dr. Álvaro Torres}

\noindent \textbf{Protocol:} IML-BA 1979.06.12-1260
\textbf{Patient:} Joaquim Antônio Correia (Hospital Record 1,260)
\textbf{Date of Death:} June 12, 1979
\textbf{Forensic Pathologist:} Dr. Álvaro Torres, CRM-MG 3,117

\vspace{0.5em}\noindent\rule{\textwidth}{0.5pt}\vspace{0.5em}

\noindent \textbf{Forensic Medical Institute of Barbacena --- Autopsy Report No. 1,260 --- June 14, 1979}

\noindent \textbf{Patient:} Joaquim Antônio Correia
\textbf{Record:} 1,260
\textbf{Estimated age:} 72 years (± 5 years)
\textbf{Place of birth:} Quixeramobim, Ceará
\textbf{Date of death:} June 12, 1979
\textbf{Date of autopsy:} June 14, 1979
\textbf{Forensic pathologist:} Dr. Álvaro Torres --- CRM-MG 3,117
\textbf{Requesting party:} Administration of the Barbacena Hospital-Colony

\vspace{0.5em}\noindent\rule{\textwidth}{0.5pt}\vspace{0.5em}

\subsection*{1. EXTERNAL EXAMINATION}

\subsubsection*{1.1 Anthropometric Data}

\noindent - \textbf{Weight:} 42 kg (body mass index: 16.0 kg/m² --- severe chronic malnutrition)
- \textbf{Height:} 1.62 m
- \textbf{Nutritional status:} Marked cachexia. Absent adipose panniculus. Globally reduced muscle mass. The bony prominences --- especially the clavicles, iliac crests and patellae --- project beneath the skin, with pronounced bony prominence. The skin is dry, with persistent skinfold (sign of chronic dehydration). The nails show Beau's lines (transverse grooves indicating repeated episodes of acute malnutrition throughout life).

\subsubsection*{1.2 Ocular Findings}

\noindent \textbf{Unilateral asymmetric jaundice --- right eyeball.}

\noindent The sclera of the right eye presents an intense, unusual golden-yellow discoloration, extending from the cornea to the bulbar conjunctiva in a homogeneous tone. The coloration is not compatible with classic hepatic jaundice (which would be bilateral and diffuse) nor with palpebral xanthelasmas. The sclera of the left eye is white, with normal vascularization.

\noindent \textit{Pathologist's note:} On palpation, the right eyeball presents a slightly firmer consistency than the left. There are no signs of recent trauma. There is no prosthesis. The yellow sclera does not fade under digital compression.

\noindent \textit{Hypotheses:}
- Sequela of untreated ocular trauma
- Undocumented neonatal infection
- \textit{(Marginal annotation, different handwriting: ``After the exhumation, this hypothesis fell apart. The eye glowed in the dark. There is no medical explanation for the phenomenon. --- Lúcia M.'')}

\subsubsection*{1.3 Electroconvulsive Therapy (ECT) Injuries}

\noindent The skull presents \textbf{bilateral electric burn marks} on the temporal regions, immediately anterior to the tragus of both ears. The marks are oval, approximately 3 cm × 2 cm, brownish in color, with slightly raised borders and thick scar tissue. The skin in these regions presents permanent alopecia --- the hair follicles were destroyed by the repeated passage of electric current.

\noindent \textit{Analysis:} The marks are compatible with repeated electroshock application without conductive gel, using dry-contact metal electrodes. The absence of conductive gel --- standard practice at the Colony until the 1960s --- produced partial-thickness burns which, upon healing, destroyed the sweat glands and hair follicles of the region.

\noindent \textbf{ECT-compatible fractures:}
- \textbf{Compression fracture of the T5 vertebra} (fifth thoracic vertebra), with loss of approximately 30\% of the vertebral body height. Old healing, without exuberant callus. Typical of violent, uncontrolled muscle contraction during a generalized tonic-clonic seizure --- the patient was lying supine during the procedure, without adequate restraint, and the forced flexion of the spine during the tonic phase of the seizure compressed the vertebra until it fractured.
- \textbf{Consolidated fracture of the left femoral neck}, with shortening of approximately 2 cm of the left lower limb (observed during ambulation, according to nursing records). The fracture occurred during an electroshock session in 1948 --- the patient, in the clonic phase of the seizure, was thrown off the stretcher and landed on his hip.
- \textbf{Consolidated fracture of the right mandibular ramus}, with slight facial asymmetry. During the seizure, the masseter muscles contracted with enough force to fracture the jaw. The absence of oral restraint (bite guards or rubber rolls) contributed to the injury.

\subsubsection*{1.4 Transorbital Lobotomy Marks (Freeman Procedure)}

\noindent \textbf{Both inner corners of the eyes} present punctate scars, approximately 2 mm in diameter, located at the junction between the lacrimal caruncle and the semilunar fold. The scars are bilateral and symmetrical.

\noindent \textit{Analysis:} The scars are compatible with the introduction of a transorbital leucotome --- an instrument similar to an ice pick, approximately 15 cm long and 6 mm in diameter, with a blunt tip and graduated shaft. The procedure, described by Walter Freeman in 1946, consisted of:

\noindent 1. Subjecting the patient to electroshock to induce unconsciousness (``convulsion anesthesia'')
2. Positioning the leucotome over the ocular conjunctiva, at the inner corner of the eye
3. Tilting the instrument 20 degrees relative to the midline
4. Tapping the leucotome with a surgical mallet until it crossed the thin bony plate of the orbit (approximately 2--3 mm thick)
5. Advancing the instrument 5--7 cm into the skull
6. Moving the leucotome handle in a fanning motion (lateromedial and mediolateral) to sever the association fibers of the prefrontal lobe
7. Repeating on the contralateral side

\noindent \textit{Relevant finding:} In the cribriform plate (roof of the orbit), two small holes of approximately 4 mm in diameter are identified, with perfectly smooth, healed borders, positioned symmetrically 2.5 cm from the midline. The distance between the holes and the angulation of the bony margins confirm that the instrument was introduced toward the medial frontal lobe. \textbf{The holes are patent} --- there was no complete bony healing.

\noindent \textit{Pathologist's marginal annotation:} ``Despite the evident holes in the bone, macroscopic examination of the brain reveals no cortical scars compatible with the procedure. The lobotomy may not have been completed, or the instrument may have followed an anatomical path that did not damage the brain tissue detectably. Note: the medical record contains no authorization for lobotomy. The procedure was performed without formal registration.''

\subsubsection*{1.5 Physical Restraint Marks}

\noindent \textbf{Wrists and ankles:} Bands of circumferential hyperpigmentation on both wrists (3 cm wide) and both ankles (2.5 cm wide). The skin in these areas presents a rough texture, with loss of elasticity and absence of hair. The bands are continuous and uniform, characteristic of \textbf{prolonged pressure from leather restraints or straitjackets}.

\noindent \textit{Observation:} The bands on the wrists are deep enough to have caused bilateral compression of the median nerve (carpal tunnel syndrome from restraint). The contraction of the fingers in semiflexion (claw hand) suggests that the patient remained restrained for periods exceeding 30 consecutive days on multiple occasions.

\noindent \textbf{Chest and abdomen:} Horizontal linear excoriations at nipple height, bilateral, compatible with repeated friction from leather or canvas straps (straitjacket). The excoriations present varying degrees of healing, indicating recurring use over decades.

\noindent \textbf{Hips and thighs:} Areas of alopecia and dermal fibrosis on the anterior faces of both thighs, at inguinal fold height. Compatible with prolonged bed immobilization, with the patient lying supine and the legs restrained by leather straps or wet sheets (which, upon drying, tightened against the skin, causing tissue ischemia).

\subsubsection*{1.6 Self-Harm Scars}

\noindent Multiple linear scars on both forearms, anterior and posterior faces. Symmetrical distribution, varied temporality (from recent scars, less than 6 months old, to old scars, decades old). The pattern is consistent with deep self-inflicted scratches --- the patient used his own fingernails to carve grooves into the skin, open and repeated, until reaching the subcutaneous tissue.

\noindent \textit{Differentiation:} The scars have irregular borders and variable width, characteristic of scratches (fingernails), not cuts (blade). There are no scars suggestive of a suicide attempt by venous incision. The pattern is of \textbf{compulsive, non-suicidal self-harm}, probably associated with episodes of psychomotor agitation.

\noindent \textbf{Additional observation:} The oldest scars present overlapping tattooing from soil impregnation (dirt particles incorporated into the dermis during healing), suggesting that the lesions were opened and exposed to contaminated soil. \textbf{There is no record of tetanus prophylaxis in the patient's medical record.}

\subsubsection*{1.7 Dentition and Oral Cavity}

\noindent \textbf{Absence of 12 teeth:} Upper and lower incisors (8), right lower first molar, left upper second premolar, left lower canine, right lower third molar. The remaining teeth present:

\noindent - Marked wear of the occlusal surface (friction from bruxism, common in patients subjected to ECT without a bite guard)
- Generalized caries with pulp exposure in 4 teeth
- Untreated dental fracture in 2 teeth (compatible with mandibular contraction during ECT)

\noindent \textbf{Tongue:} Two large linear scars on the longitudinal lateral borders, approximately 3 cm each. These are scars from repeated biting during seizures. The absence of a bite guard during electroshock caused the patient to bite his own tongue on multiple occasions, with enough force to traverse the muscle tissue. The scars are deep and keloidal --- the healing tissue formed palpable nodules on the edge of the tongue.

\subsubsection*{1.8 Traumatic Amputation}

\noindent \textbf{Partial amputation of the distal phalanx of the little finger of the left hand.} Old healing, with a blunt stump and absence of the ungual phalanx. There is no record of trauma justifying the amputation in the medical record documents.

\noindent \textit{Marginal annotation:} ``Healing by second intention, without suturing. Possible bite or traumatic avulsion. The patient also tore out his own nails with his teeth (record of Apr/1937). The phalanx may have been torn off by the patient himself during an episode of self-harm.''

\subsubsection*{1.9 Thermal Burns}

\noindent On the anterior face of the chest and abdomen, \textbf{scattered hypopigmentation patches} are observed, with irregular borders, approximately 2 to 8 cm in diameter. They are scars of superficial second-degree burns, already epithelialized.

\noindent \textit{Analysis:} The thermal pattern is compatible with contact with hot water or heated surfaces. The Colony's nursing records mention the practice of ``therapeutic baths'' --- patients were immersed in hot water (40--50°C) or cold water (10--15°C) for prolonged periods, as a form of ``sedation'' or ``stimulation''. The absence of a geometric pattern in the burns suggests scalding by immersion, not contact with a heated object.

\subsection*{2. INTERNAL EXAMINATION}

\subsubsection*{2.1 Cause of Death}

\noindent \textbf{Cardiorespiratory arrest} secondary to:
- Congestive heart failure
- Severe chronic malnutrition (cardiac cachexia)
- Bronchoaspiration (food aspirated into the lower airways --- compatible with dysphagia from neurological sequelae)

\subsubsection*{2.2 Cardiovascular System}

\noindent Heart with reduced weight (220 g --- expected for the biotype: 280--340 g). Brown atrophy of the myocardium (lipofuscin pigmentation, typical of prolonged malnutrition). The cardiac fibers present cytoplasmic vacuolization and loss of transverse striation. Aorta with mild atherosclerosis, incompatible with the estimated age (chronic malnutrition delayed the atherosclerotic process).

\subsubsection*{2.3 Digestive System}

\noindent \textbf{Stomach:} Empty. Absolute absence of food content. The gastric walls are thinned (average thickness of 2 mm), with atrophic mucosa and reduced rugae. The patient had not eaten for at least 24 hours before death.

\noindent \textbf{Intestines:} Signs of chronic malabsorption --- the intestinal villi are flattened and shortened (subtotal villous atrophy). Presence of \textit{Ascaris lumbricoides} (roundworms) in high parasitic load (17 specimens counted in the small intestine, 4 in the ascending colon). The parasites, each 15--25 cm long, formed a tangle occupying approximately 40\% of the intestinal lumen.

\noindent \textit{Observation:} The high parasitic load, associated with chronic malnutrition, contributed to the state of cachexia. The patient competed with his own parasites for the few nutrients he managed to ingest.

\noindent \textbf{Liver:} Reduced weight (980 g), with moderate steatosis and signs of protein-calorie malnutrition. No signs of cirrhosis or liver disease that would justify the unilateral jaundice observed in the external examination.

\subsubsection*{2.4 Central Nervous System}

\noindent \textbf{Brain:} Weight 1,120 g (within expected parameters for age and biotype). The cerebral hemispheres are symmetrical. The meninges are transparent, without signs of meningitis or subarachnoid hemorrhage.

\noindent \textit{Specific analysis --- lobotomy:} \textbf{Absence of macroscopic cortical scars} compatible with transorbital or prefrontal lobotomy. The orbitofrontal and medial prefrontal areas were dissected and carefully examined. There is no gliosis, cavitation, or parenchymal alteration indicating severing of nerve fibers. The access route of the leucotomes (through the roof of the orbits) was identified, but the trajectory appears to have been interrupted before reaching the cerebral parenchyma, or the instrument may have followed an anatomical plane that did not damage detectable brain tissue.

\noindent \textit{Pathologist's note:} ``Considering the absence of brain damage, it is possible that the lobotomy procedure was interrupted before parenchymal injury, or that the patient was subjected only to the bony perforation (as ``training'' for residents) without severing the frontal fibers. There is no record in the medical file to clarify.''

\noindent \textit{Specific analysis --- electroshock:} ``There is no macroscopic evidence that the repeated electroshock procedures (estimated at hundreds of sessions over 50 years of internment) caused detectable brain damage. Absence of gliosis, lobar atrophy, or significant vascular alterations.''

\noindent \textit{Marginal annotation in cursive handwriting, seemingly more recent:} ``The absence of brain damage is not a relief. It is the most disturbing of all findings. It means Joaquim was not violent because of brain injury. The violence was his. Or it was something else's. The procedures left no mark on the brain --- but they left marks all over the body. The marks of the body tell a story. The silent brain is the loudest part of the report. --- L.M.''

\subsubsection*{2.5 Vertebral Column}

\noindent \textbf{Compression fracture at T5} (as described in item 1.3). The spinal cord at the T5--T6 level presents slight anteroposterior flattening, without significant compression. The fracture was of the ``stable'' type and caused no neurological deficit. \textbf{Additional compression fractures at T3 and T8}, smaller, not identified in the previous radiological examination.

\noindent \textit{Observation:} The three fractured vertebrae (T3, T5, T8) form a classic triad of ECT fractures --- they occur during the tonic phase of the seizure, when the simultaneous contraction of the flexor and extensor muscles of the spine generates axial compressive forces that exceed bone resistance.

\subsubsection*{2.6 Musculoskeletal System --- Additional Findings}

\noindent - \textbf{Shoulder arthropathy:} Recurrent bilateral glenohumeral dislocation (shoulders dislocated and reduced on multiple occasions). Associated with the restraint maneuver known as the ``arm lock'' --- a technique used by the Colony's orderlies to immobilize agitated patients, which consists of hyperabducting and externally rotating the arm until the joint pops.
- \textbf{Scalp scar:} Linear 6 cm scar on the occipital region, compatible with a blunt-force wound. Possible fall during a seizure or physical assault by staff.
- \textbf{Foot deformity:} Severe bilateral hallux valgus, with overlapping of the second toe over the hallux on both feet. Deformity caused by prolonged use of inadequate footwear (one-size, unnumbered --- ``Colony slippers'', made by the patients themselves in the hospital's shoemaking workshop).

\subsection*{3. SUMMARY OF PROCEDURES IDENTIFIED ON THE BODY}

\begin{moltablethree}
\hline
\textbf{Procedure} & \textbf{Physical Evidence} & \textbf{Period} \\
\hline
Electroconvulsive therapy without anesthesia & Temporal burns; vertebral fractures T3/T5/T8; mandibular fracture; femur fracture; lingual scars; tooth loss & 1925--1979 (54 years) \\
Transorbital lobotomy (Freeman) & Holes in the cribriform plate, bilateral; conjunctival scars & 1950s (unrecorded) \\
Straitjacket / prolonged restraint & Hyperpigmentation bands on wrists and ankles; digital contracture; pressure alopecia & 1920--1979 (59 years) \\
Therapeutic baths (scalding) & Thermal burn scars on chest and abdomen & 1930--1960 \\
Insulin therapy (insulin coma) & \textit{(Absence of marker --- leaves no physical scars)} & 1940--1950 \\
Restraint with wet sheets & Dermal fibrosis on thighs and hips & 1930--1979 \\
\hline
\end{moltablethree}

\noindent \textit{Final note (manuscript, added later):} The body presents 59 years of markers of institutional intervention: ECT without anesthesia (bilateral temporal burns, vertebral fractures T3/T5/T8), transorbital lobotomy (bilateral perforation of the cribriform plate), mechanical restraint (digital contracture, dermal fibrosis on the thighs), therapeutic baths (thermal scars on chest and abdomen). \textbf{Cause of death:} cardiorespiratory arrest. \textbf{Contributing cause:} cachexia from chronic malnutrition (estimated weight: 38 kg). \textbf{Comorbidity:} high parasitic load (\textit{Ascaris lumbricoides}, 21 specimens). No markers of organic disease were found that would justify 62 years of institutionalization. The patient's most significant injury is not anatomical.

\vspace{0.5em}\noindent\rule{\textwidth}{0.5pt}\vspace{0.5em}

\noindent Dr. Álvaro Torres

\noindent Forensic Pathologist --- IML Barbacena

\noindent CRM-MG 3,117

\vspace{0.5em}\noindent\rule{\textwidth}{0.5pt}\vspace{0.5em}

\noindent \textit{Addendum (manuscript, seemingly added later):} ``This report was written on June 14, 1979, before the exhumation. The findings described in the following section --- Exhumation and Special Findings --- were added after the opening of the grave and the recovery of the diary. The signatory declares that there is no medical explanation for the phenomena described in item 4 (Exhumation), limiting himself to reporting what was presented by the gravedigger and by the recovered material. The diary remains in the custody of the IML, attached to this case file. --- A.T.''

\noindent \textit{Additional note (L.M.'s handwriting, 2026 protocol):} ``Bone fragments of Strigiformes (burrowing owl) were recovered from the same stratum of the grave, alongside the patient's material. The presence of owl remains in Joaquim's burial was recorded by the gravedigger as 'the bird that would not leave'. No taphonomic explanation was found for the deposition. Owl bones are not a common finding in human burials at the Colony.''

\noindent\textit{\textbf{↪ Routes:} \rota{MEM-04} \rotasep \rota{DOC-02} \rotasep \rota{LUC-03} \rotasep \rota{CONT-03}}
